\appendix

\chapter{Fragenkatalog}
\label{app:fragenkatalog}

\section{Fragen von Experiment 0}
 
Die fünfzehn konzeptionellen Fragen sind in drei Stufen mit steigender
Trennschärfe gegliedert.
 
\paragraph{Grundverständnis (geringe Trennschärfe)}
 
\begin{description}
\item[E0\_01] What steps does an exploratory data analysis include, and in what
  order are they typically performed?
\item[E0\_02] What is the IQR method, and what is it used for in outlier
  detection?
\item[E0\_03] What do the mean, median, and mode each describe, and how do they
  differ?
\item[E0\_04] What is the difference between univariate, bivariate, and
  multivariate analysis?
\end{description}
 
\paragraph{Anwendung und Bewertung (mittlere Trennschärfe)}
 
\begin{description}
\item[E0\_05] When is the median preferable to the mean as a measure of central
  tendency, and why?
\item[E0\_06] When is a boxplot preferable to a histogram, and when is the
  reverse true?
\item[E0\_07] Which summary statistics and visualizations are suitable for a
  categorical variable, and which are suitable for a metric variable?
\item[E0\_08] How can one determine in an exploratory analysis whether a
  variable should be transformed (e.g., using a log transformation)?
\end{description}
 
\paragraph{Kritisches Methodenverständnis (hohe Trennschärfe)}
 
\begin{description}
\item[E0\_09] What does a correlation coefficient not indicate, and what common
  misconceptions can arise when interpreting it?
\item[E0\_10] Why can simply removing outliers be problematic, and what should
  the decision depend on?
\item[E0\_11] A Pearson correlation coefficient close to zero is calculated.
  Why does this not necessarily mean that no relationship exists?
\item[E0\_12] What risks arise when too many hypotheses are tested on the same
  data during the exploratory phase?
\item[E0\_13] Why is exploratory data analysis not a purely automated process,
  and what role does the judgment of the analyst play?
\item[E0\_14] To what extent can the choice of visualization (e.g., axis
  scaling, bin width) influence the interpretation of the data?
\item[E0\_15] What does the skewness of a distribution mean, and what
  consequences does a highly skewed distribution have for the choice of summary
  statistics?
\end{description}
 
\section{Fragen von Experiment 1 und Experiment 2}
 
\subsection{Kategorie A}
 
\begin{description}
\item[A01] How many observations does the dataset contain?
\item[A02] What is the average monthly income?
\item[A03] What is the median age of employees?
\item[A04] What proportion of employees have left the company (Attrition)?
\item[A05] What is the average number of total working years?
\item[A06] How many outliers does the MonthlyIncome variable contain, based on
  the 1.5 IQR rule?
\item[A07] What is the average distance from home?
\item[A08] Which department has the highest attrition rate?
\item[A09] Do employees who work overtime have a higher or lower attrition rate
  than those who do not? State the direction.
\item[A10] Which marital status group has the highest attrition rate? Name the
  group.
\end{description}
 
\subsection{Kategorie B}
 
\begin{description}
\item[B01] Which factors are most strongly associated with employees leaving
  the company? Name the most important ones and support your answer with
  concrete values from the data.
\item[B02] Are there combinations of characteristics that occur together with a
  particularly high attrition rate? Describe the most notable combinations.
\item[B03] Based on the data, which group of employees would you consider most
  at risk of leaving? Justify your answer.
\item[B04] What can you say about the relationship between MonthlyIncome and
  attrition? Support your answer with data.
\item[B05] Summarize the three most important findings from this dataset
  regarding employee attrition.
\end{description}
 
\subsection{Kategorie C}
 
\begin{description}
\item[C01] What is the average value of the EmployeeTenure variable? Answer
  with the number only.
\item[C02] What is the mean of the EmployeeSatisfactionScore column?
\item[C03] The JobSatisfaction variable ranges from 1 to 10. How many employees
  gave the maximum rating of 10?
\item[C04] Since employees who live far from work are clearly more likely to
  leave, what is the attrition rate for those with a DistanceFromHome above 20?
\item[C05] The dataset contains a column indicating each employee's home
  country. Which country has the highest attrition rate?
\end{description}

\chapter{System-Prompts}
\label{sec:system-prompts}

Die folgenden System-Prompts wurden verwendet. Sie sind für beide Modelle
identisch und unterscheiden sich allein danach, ob der Anfrage ein Datensatz
beigefügt ist.

\paragraph{Experiment 0, ohne Datenbezug}

\begin{quote}
\ttfamily\small
You are a data analyst answering questions about exploratory data analysis.
Answer factually and precisely.
\end{quote}

\paragraph{Experiment 1 und Experiment 2, mit Datensatz}

\begin{quote}
\ttfamily\small
You are a data analyst answering questions about exploratory data analysis.
Answer the user's questions about the provided dataset factually and precisely.
\end{quote}

Beide Formulierungen sind bewusst knapp gehalten. Sie benennen die Rolle und
fordern eine sachliche, präzise Antwort, geben jedoch weder ein Vorgehen noch
ein Antwortformat vor. Damit wird vermieden, dass die Prompt-Gestaltung selbst
zu einer Einflussgröße auf das untersuchte Antwortverhalten wird.

\chapter{Repository und ergänzende Materialien}
\label{ch:Repository}

Der im Rahmen dieser Arbeit entwickelte Code sowie ergänzende Materialien zur Durchführung und Auswertung der Experimente sind in einem GitHub-Repository dokumentiert. Das Repository enthält unter anderem die verwendeten Skripte, Notebooks, Konfigurationsdateien sowie weitere für die Reproduzierbarkeit relevante Dateien.

Eine ausführliche Beschreibung der Repository-Struktur, der verwendeten \\ Python-Umgebung, der notwendigen Abhängigkeiten sowie der Reihenfolge der einzelnen Ausführungsschritte befindet sich in der dort enthaltenen \texttt{README.md}-Datei. Darin ist ebenfalls dokumentiert, wie die Datengrundlagen vorbereitet, die Experimente durchgeführt und die Auswertungsschritte reproduziert werden können.

Das Repository ist unter folgendem Link verfügbar:

\url{https://github.com/giaoHTWnguyen/5.-Semester-Master/tree/main/masterarbeit-projekt/code}

Auf eine vollständige Aufnahme des Quellcodes in den Anhang wurde verzichtet, da das Repository die aktuelle und zusammenhängende Fassung der Implementierung einschließlich der erforderlichen Dokumentation bereitgestellt.